\section{Introduction}

Misinformation and fake news have become a growing threat in today's digital landscape, spreading rapidly across social media platforms and influencing public opinion. Unlike traditional text-based misinformation, modern fake news often combines both textual and visual content, making unimodal detection approaches insufficient for capturing the full context of deceptive information.

Fakeddit, introduced by Nakamura \cite{fakeddit}, is a large-scale multimodal fake news detection dataset collected from Reddit. It contains over one million samples with both text and image modalities, annotated with fine-grained labels for 2-way and 6-way classification tasks. The original paper proposed multimodal baseline models that combine text and image features, demonstrating the importance of leveraging both modalities for fake news detection.

In this work, we reproduce the baseline models proposed in the Fakeddit paper, including text-only, image-only, and multimodal variants using a subset of data which contains 50.000 samples.

Building upon this reproduction, we propose an extension to the multimodal model by incorporating a cross-attention mechanism \cite{attention} between the text and image representations. Unlike simple feature combination used in the original baseline, cross-attention allows the model to dynamically align and weight relevant information across modalities, enabling richer and more expressive multimodal fusion.